\chapter{Complément : la géométrie des probabilités}
\section*{Exercice 28.1 -- Moyennes de groupes}
La moyenne du premier groupe est $2$, celle du second $6$. La variable conditionnelle vaut donc $Y=(2,2,6,6)$. Sa moyenne est $4$, égale à celle de $X$ : $(1+3+4+8)/4=4$.

\section*{Exercice 28.2 -- Variance totale}
La variance totale est $[(1-4)^2+(3-4)^2+(4-4)^2+(8-4)^2]/4=26/4=6{,}5$. Les variances internes valent $1$ et $4$, de moyenne $2{,}5$. La variance de $Y=(2,2,6,6)$ vaut $4$. On retrouve $6{,}5=2{,}5+4$.

\section*{Exercice 28.3 -- Erreur d'un prédicteur}
L'erreur de $Y$ vaut $[(1)^2+(1)^2+(-2)^2+2^2]/4=2{,}5$. Comme $Y-Z$ vaut constamment $1$, son carré moyen vaut $1$. L'identité de projection donne une erreur de $Z$ égale à $3{,}5$. Directement, les carrés des résidus sont $0,4,1,9$, de moyenne $3{,}5$.

\section*{Exercice 28.4 -- Corrélation extrême}
La covariance vaut $-3\operatorname{Var}(X)$ et $\sigma(Y)=3\sigma(X)$. Donc $\rho=-3\sigma(X)^2/[3\sigma(X)^2]=-1$. Il s'agit d'une relation affine décroissante exacte, presque sûrement.

\section*{Exercice 28.5 -- Jensen}
La dérivée de l'exponentielle est croissante, donc sa courbe est au-dessus de la tangente en $\mu=\mathbb E(X)$ : $e^x\ge e^\mu+e^\mu(x-\mu)$. Prendre l'espérance annule le terme centré et donne l'inégalité.

\section*{Exercice 28.6 -- Information complète}
Dans le premier cas, chaque moyenne de groupe est simplement la valeur de $X$ sur l'issue : $Y=X$ et l'erreur de prédiction est nulle. Dans le second, le seul groupe a pour moyenne $\mathbb E(X)$ : le prédicteur est constant et l'erreur quadratique minimale est la variance.
